\section{Introducción}
    Una recta que explica el 99,7\% de la variabilidad de los datos parece un
    resultado difícil de mejorar. Este informe muestra que ese número, por sí solo,
    no dice si el modelo es correcto, dado que la misma regresión que alcanza un
    coeficiente de determinación de 0,9969 deja residuos que no son ruido, sino que
    se ordenan según una variable que quedó fuera del modelo. Detectar esa
    estructura, incorporarla y verificar que el modelo corregido cumple los
    supuestos que sostienen su inferencia es el recorrido completo de un análisis de
    regresión \cite{montgomery2021regression}, \cite{wooldridge2020econometrics}, y
    es también el recorrido de este trabajo.

    \vspace{10pt}

    El análisis se apoya en un conjunto de datos simulado del consumo energético
    mensual de 120 clientes de una empresa distribuidora, repartidos en los sectores
    Residencial, Comercial e Industrial. Que el conjunto sea simulado es una ventaja
    metodológica y no una limitación, puesto que se sabe de antemano que existen
    tres poblaciones con tarifas medias distintas, de modo que el modelo puede
    evaluarse por su capacidad de recuperar una estructura conocida y no solo por su
    ajuste.

    \vspace{10pt}

    El informe se organiza en dos capítulos. El primero describe la metodología,
    con el conjunto de datos, la secuencia de modelos estimados, las medidas que los
    comparan y el flujo de trabajo reproducible que los produce; el segundo presenta
    los resultados, desde el análisis de correlación hasta la verificación cruzada
    entre Python y R, pasando por el diagnóstico del modelo simple, la especificación
    múltiple, la validación predictiva y la visualización avanzada
    \cite{castro2026avanzada}. Los seis scripts que sostienen el análisis se
    reproducen completos en los Apéndices \ref{anexo_python} y \ref{anexo_r}, de modo
    que cada cifra del informe puede rastrearse hasta la línea que la produce.

    \vspace{10pt}

    Tres hallazgos anticipan el contenido. El primero es que las pendientes del
    modelo múltiple tienen nombre propio, pues son las tarifas de cada sector,
    791,6, 704,8 y 671,0 pesos por kilovatio hora, y reproducen con menos del 4\% de
    diferencia el cociente entre costo y consumo calculado sin pasar por la
    regresión. El segundo es que ese modelo no solo ajusta mejor dentro de la
    muestra, sino que también predice mejor fuera de ella, ya que la validación
    cruzada de diez pliegues reduce el error un 16,7\% frente al modelo simple. El
    tercero es que el sector Industrial paga 121 pesos por kilovatio hora menos que
    el Residencial, un descuento por escala que el análisis cuantifica y que la
    empresa puede contrastar con su política tarifaria.
